\section{Reaction networks}

The reaction-network fiber records the counted, consumable regime above
monotone closure.  Its organizing projection is not a new primitive: a reaction
state is a finite resource state, firing is the nonmonotone counted step,
Datalog is the support projection obtained by forgetting counts, and the
thermodynamic/rate results are projections of that same carrier.

\begin{proposition}[Reaction networks are counted resource systems]\label{thm:reaction-network-carrier}\lean{Infomaterial.Relational.Reaction.Network}\leanok\uses{thm:closure-rule-lfp, thm:lambek-biresiduation}
A reaction network is a finite roster of counted states and consumable
reactions.  It sits above monotone closure: counts can be spent, so the
resource carrier is not itself an inflationary closure system.
\end{proposition}

\begin{theorem}[Balanced firing conserves mass]\label{thm:reaction-mass-conserved}\lean{Infomaterial.Relational.Reaction.mass_conserved}\leanok\uses{thm:reaction-network-carrier}
For a balanced reaction, one enabled firing preserves the total count over the
chosen finite roster.
\end{theorem}

\begin{theorem}[Reaction firing is not monotone]\label{thm:reaction-fire-not-monotone}\lean{Infomaterial.Relational.Reaction.fire_not_monotone}\leanok\uses{thm:reaction-network-carrier, thm:closure-rule-lfp}
Reaction firing can consume a resource.  This is the formal obstruction to
identifying counted reaction dynamics with an inflationary Datalog closure.
\end{theorem}

\begin{theorem}[Catalysis is fixed under firing]\label{thm:reaction-catalyst-fixed}\lean{Infomaterial.Relational.Reaction.catalyst_is_fixed}\leanok\uses{thm:reaction-network-carrier}
A catalyst that appears with the same count on the reactant and product sides
is unchanged by the firing it enables.
\end{theorem}

\begin{theorem}[Datalog is the affine support projection]\label{thm:reaction-support-projection}\lean{Infomaterial.Relational.Reaction.support_fire_le_datalog}\leanok\uses{thm:reaction-fire-not-monotone, thm:closure-rule-lfp}
After forgetting counts to supports, an enabled reaction lands below the
monotone Datalog immediate-consequence step.  The missing support is exactly
the consumption the affine support projection cannot record.
\end{theorem}

\begin{theorem}[Reachability is the settled trajectory relation]\label{thm:reaction-reachability}\lean{Infomaterial.Relational.Reaction.reach_trans}\leanok\uses{thm:reaction-network-carrier}
The reaction reachability relation is closed under composition of firing
steps, giving the incidence carrier for trajectories through the counted
resource system.
\end{theorem}

\begin{theorem}[Mass is invariant along balanced reachability]\label{thm:reaction-reach-conserves-mass}\lean{Infomaterial.Relational.Reaction.reach_conserves_mass}\leanok\uses{thm:reaction-reachability, thm:reaction-mass-conserved}
If every reaction in the network is balanced, total mass is preserved along
every reachable trajectory.
\end{theorem}

\begin{theorem}[Catabolic trajectories descend in free energy]\label{thm:reaction-catabolic-descent}\lean{Infomaterial.Relational.Reaction.catabolic_descent}\leanok\uses{thm:reaction-reachability, thm:reaction-fire-not-monotone}
For a catabolic network, step-counted reachability carries a natural-valued
Lyapunov bound: each firing spends at least one unit of total count.
\end{theorem}

\begin{theorem}[Mass-action propensities are nonnegative]\label{thm:reaction-propensity-nonnegative}\lean{Infomaterial.Relational.Reaction.propensity_nonneg}\leanok\uses{thm:euler-e-internal-real, thm:reaction-network-carrier}
With nonnegative concentrations and rate constants, a mass-action propensity is
nonnegative in the internal real carrier.
\end{theorem}

\begin{theorem}[Equilibria are Euler fixed points]\label{thm:reaction-mass-action-equilibrium}\lean{Infomaterial.Relational.Reaction.equilibrium_is_fixed}\leanok\uses{thm:reaction-propensity-nonnegative, thm:real-stochastic-interval-value-category}
A concentration at which the rate field vanishes is fixed by every discrete
Euler step of that field.
\end{theorem}

\begin{theorem}[Detailed balance is zero net flux]\label{thm:reaction-detailed-balance}\lean{Infomaterial.Relational.Reaction.detailedBalance_iff_netFlux_zero}\leanok\uses{thm:reaction-mass-action-equilibrium}
For a reversible reaction pair, equality of forward and backward propensities
is equivalent to zero net flux.
\end{theorem}

\begin{theorem}[Reaction entropy relaxes under balancing]\label{thm:reaction-entropy-relaxes}\lean{Infomaterial.Relational.Reaction.entropyH_le_of_relaxes}\leanok\uses{thm:finite-boltzmann-collision-lfp, thm:reaction-reach-conserves-mass}
When the occupation multiset of a firing trajectory relaxes by finite
Boltzmann exchanges, the reaction-state \(H\) functional does not increase.
\end{theorem}

\begin{theorem}[Continuum reaction chemistry frontier]\label{thm:reaction-continuum-frontier}\notready\uses{thm:reaction-support-projection, thm:reaction-catabolic-descent, thm:reaction-detailed-balance, thm:reaction-entropy-relaxes}
The formalized reaction sector is still finite and discrete.  The continuum
reaction-chemistry frontier needs measure, stochastic-process, smooth ODE/PDE,
and analytic convergence layers before these finite projections become a full
chemical kinetics theory.
\end{theorem}
